First, a necessary correction.

The literal claim in part \( (1) \),
\[
|\operatorname{vert}(P+Q)|\ge 4 \quad\text{for every triangle } P \text{ and every } |T|\ge 2,
\]
is **false as stated**. For example, take
\[
S=T=\{(0,0),(1,0),(0,1)\},
\]
and choose \(\varepsilon\) with both signs, say
\[
\varepsilon(0,0)=+1,\qquad \varepsilon(1,0)=+1,\qquad \varepsilon(0,1)=-1.
\]
Then \(P=Q\) is the standard triangle, so
\[
P+Q=2P
\]
is again a triangle, hence
\[
|\operatorname{vert}(P+Q)|=3.
\]
So that geometric sentence cannot be proved.

Likewise, the geometric hint for part \( (2) \) is not sufficient by itself: if \(P\) is any non-parallelogram convex quadrilateral and \(Q=P\), then
\[
P+Q=2P
\]
still has \(4\) vertices.

So what I will prove is the **correct extremal conclusion**:

\[
n=3 \implies |\operatorname{supp}(F)|\neq 3,
\]
and
\[
n=4,\ |\operatorname{supp}(F)|=4 \implies S \text{ is the vertex set of a parallelogram}.
\]

---

## \(1.\) Triangle case: \( |\operatorname{supp}(F)|=3 \) is impossible

Let
\[
L:=\langle S-S\rangle_{\mathbb Z}.
\]
Since \(S\) is a triangle, \(L\) has rank \(2\).

Decompose \(T\) into cosets of \(L\):
\[
T=\bigsqcup_\alpha T_\alpha.
\]
For each \(\alpha\), define
\[
F_\alpha(x):=\sum_{t\in T_\alpha}\varepsilon(t)\mathbf 1_{S+t}(x).
\]
Because every set \(S+t\) lies in the coset \(t+L\), the supports of the \(F_\alpha\) are pairwise disjoint. Hence
\[
F=\sum_\alpha F_\alpha
\quad\text{and}\quad
\operatorname{supp}(F)=\bigsqcup_\alpha \operatorname{supp}(F_\alpha).
\]
If \(T_\alpha\neq\varnothing\), then the same lonely-point/vertex argument as in Subproblem \(1\) gives
\[
|\operatorname{supp}(F_\alpha)|\ge |\operatorname{vert}(P+\operatorname{conv}(T_\alpha))|\ge 3.
\]
Therefore, if \( |\operatorname{supp}(F)|=3 \), there can be only **one** nonempty coset. After translating, we may thus assume
\[
T\subset L.
\]

Now choose a \(\mathbb Z\)-basis \(u,v\) of \(L\) so that, after translating one vertex of \(S\) to the origin,
\[
S=\{0,u,v\}.
\]
Encode points of \(L\) by Laurent monomials, with \(u\mapsto X\) and \(v\mapsto Y\). Then
\[
A(X,Y):=1+X+Y,
\]
and if
\[
G(X,Y):=\sum_{t\in T}\varepsilon(t)X^{a_t}Y^{b_t},
\]
then
\[
H(X,Y):=A(X,Y)G(X,Y)
\]
has coefficients exactly equal to the values of \(F\). So
\[
|\operatorname{supp}(H)|=|\operatorname{supp}(F)|.
\]

Assume for contradiction that
\[
|\operatorname{supp}(F)|=3.
\]
Then \(H\) is a trinomial divisible by \(1+X+Y\).

### Lemma
If a trinomial \(H\in \mathbb Q[X^{\pm1},Y^{\pm1}]\) is divisible by \(1+X+Y\), then \(H\) is a monomial multiple of \(1+X+Y\).

#### Proof
Factor out a monomial from \(H\); we may write
\[
H=a+bX^{p}Y^{q}+cX^{r}Y^{s},
\qquad a,b,c\neq 0.
\]
Since \(1+X+Y\mid H\), substituting \(Y=-1-X\) gives the identity
\[
a+bX^{p}(-1-X)^q+cX^{r}(-1-X)^s\equiv 0.
\]

Look at the order at \(X=0\). The first term has order \(0\), so at least one of the other two must also have order \(0\). Thus one of \(p,r\) is \(0\).

Now look at the order at \(X=-1\), i.e. at \(Z:=X+1=0\). Again the constant term has order \(0\), so one of \(q,s\) is \(0\).

Since the exponent pairs are distinct, after relabeling we get
\[
H=a+bX^m+cY^n
\qquad (m,n>0).
\]
Thus
\[
a+bX^m+c(-1-X)^n\equiv 0.
\]

Comparing highest degrees yields
\[
m=n=:d
\quad\text{and}\quad
b=-c(-1)^d.
\]
Evaluating at \(X=0\) gives
\[
a+c(-1)^d=0,
\]
so
\[
a=-c(-1)^d=b.
\]
Evaluating at \(X=-1\) gives
\[
a+b(-1)^d=0.
\]
Since \(a=b\), this becomes
\[
a(1+(-1)^d)=0,
\]
hence \(d\) is odd. Therefore
\[
a=b=c.
\]
So
\[
1+X^d+(-1-X)^d\equiv 0.
\]
But the coefficient of \(X^{d-1}\) in the left-hand side is \(-d\), so this identity is possible only for
\[
d=1.
\]
Hence \(H\) is a monomial multiple of \(1+X+Y\). ∎

Applying the lemma, \(H\) is a monomial multiple of \(A=1+X+Y\). Therefore \(G\) is a monomial.

But \(G\) was
\[
G=\sum_{t\in T}\varepsilon(t)X^{a_t}Y^{b_t},
\]
and \(T\) has both signs, so \(G\) is **not** a monomial. Contradiction.

Therefore
\[
|\operatorname{supp}(F)|\ge 4
\]
in the triangle case, and in particular equality
\[
|\operatorname{supp}(F)|=3
\]
is impossible.

---

## \(2.\) Quadrilateral case: if \( |\operatorname{supp}(F)|=4 \), then \(S\) is a parallelogram

Now assume \(n=4\), so \(S\) is the vertex set of a convex quadrilateral, and suppose
\[
|\operatorname{supp}(F)|=4.
\]

As above, encode the finite additive subgroup generated by \(S\cup T\) by Laurent monomials and define
\[
A:=\sum_{s\in S} z^s,
\qquad
G:=\sum_{t\in T}\varepsilon(t) z^t,
\qquad
H:=AG.
\]
Then
\[
|\operatorname{supp}(H)|=|\operatorname{supp}(F)|=4.
\]
So \(H\) is a \(4\)-term Laurent polynomial, and since \(T\) has both signs, \(G\) is not a monomial; hence \(H\) is reducible.

We need one elementary lemma.

### Lemma
Let
\[
A=\sum_{s\in S} z^s
\]
with \(|S|=4\) and all coefficients equal to \(1\). If
\[
B=\alpha z^u+\beta z^v
\]
is a binomial and \(AB\) has exactly \(4\) monomials, then \(S\) is a parallelogram.

#### Proof
Let
\[
d:=v-u\neq 0.
\]
Up to the monomial factor \(z^u\),
\[
AB=z^u\bigl(\alpha A+\beta z^dA\bigr).
\]
So the relevant exponent sets are \(S\) and \(S+d\).

Let
\[
I:=S\cap (S+d).
\]
At each exponent in \(I\), exactly one monomial from \(A\) and one from \(z^dA\) meet, so the coefficient there is \(\alpha+\beta\).

- If \(\alpha+\beta\neq 0\), then no intersection exponent disappears, so
  \[
  |\operatorname{supp}(AB)|=|S\cup (S+d)|=8-|I|.
  \]
  Since \(d\neq 0\), we cannot have \(S=S+d\), so \(|I|\le 3\), and hence
  \[
  |\operatorname{supp}(AB)|\ge 5,
  \]
  contradiction.

Thus necessarily
\[
\alpha+\beta=0.
\]
Then every exponent in \(I\) cancels, and
\[
|\operatorname{supp}(AB)|=8-2|I|.
\]
Since this equals \(4\), we get
\[
|I|=2.
\]
So there are two distinct pairs of points of \(S\) with the same nonzero difference \(d\). Equivalently, there exist distinct \(p_1,p_2,p_3,p_4\in S\) such that
\[
p_1-p_2=p_3-p_4\neq 0.
\]
Hence
\[
p_1+p_4=p_2+p_3.
\]
A \(4\)-point set in convex position with this relation is exactly a parallelogram. ∎

---

Now return to \(H\).

Choose a \(\mathbb Z\)-basis for the subgroup generated by the exponent vectors of the four monomials of \(H\), and multiply by a monomial so that \(H\) becomes an ordinary quadrinomial over \(\mathbb Q\).

Schinzel's classification of reducible quadrinomials says the following: over characteristic \(0\), a reducible quadrinomial with exponent-rank \(3\) cannot occur; and when the exponent-rank is \(2\), every reducible quadrinomial is either one of three explicit exceptional families or else has a binomial factor, with complementary factor irreducible or binomial. ([nzdr.ru](https://www.nzdr.ru/data/media/biblio/kolxoz/M/MA/MAco/Schinzel%20A.%20Polynomials%20with%20special%20regard%20to%20reducibility%20%28CUP%2C%202000%29%28ISBN%200521662257%29%28O%29%28569s%29_MAco_.pdf))

Since \(H\) is reducible, its exponent-rank cannot be \(3\). So its exponent-rank is \(2\), and Schinzel's rank-\(2\) classification applies. In the three exceptional families, the nonmonomial irreducible factors have either \(3\) terms, \(6\) terms, or \(4\) terms with coefficient pattern \(1,\pm1,-2,\pm2\). But \(A\) is a factor of \(H\), and \(A\) has exactly \(4\) terms, all with coefficient \(1\). So \(A\) cannot be one of those exceptional irreducible factors. Therefore \(H\) is **not** exceptional. ([nzdr.ru](https://www.nzdr.ru/data/media/biblio/kolxoz/M/MA/MAco/Schinzel%20A.%20Polynomials%20with%20special%20regard%20to%20reducibility%20%28CUP%2C%202000%29%28ISBN%200521662257%29%28O%29%28569s%29_MAco_.pdf))

Hence \(H\) falls into the non-exceptional case:
\[
H=DR,
\]
where \(D\) is a binomial and \(R\) is irreducible or binomial. ([nzdr.ru](https://www.nzdr.ru/data/media/biblio/kolxoz/M/MA/MAco/Schinzel%20A.%20Polynomials%20with%20special%20regard%20to%20reducibility%20%28CUP%2C%202000%29%28ISBN%200521662257%29%28O%29%28569s%29_MAco_.pdf))

Now \(A\mid H\). I claim that \(A\) cannot divide a binomial. Indeed, if
\[
D=\alpha z^u+\beta z^v,
\]
write
\[
v-u=md_0
\]
with \(d_0\) primitive. Over \(\overline{\mathbb Q}\),
\[
D=\alpha z^u\prod_{\zeta^m=-\beta/\alpha}(1-\zeta z^{d_0}),
\]
so every divisor of \(D\) is a monomial times a polynomial in the single monomial \(z^{d_0}\). Therefore every divisor of \(D\) has all exponent vectors on one affine line. But the support of \(A\) is the \(4\)-point set \(S\) in convex position, so it is not collinear. Thus
\[
A\nmid D.
\]
The same argument shows that \(A\) cannot divide a binomial \(R\) either.

So \(R\) must be irreducible, and since \(A\mid R\), necessarily \(A\) is associated to \(R\). Consequently \(G\) is associated to the binomial \(D\). Thus \(G\) itself is a binomial (up to a nonzero scalar):
\[
G=\alpha z^u+\beta z^v.
\]

But then
\[
H=AG
\]
has exactly \(4\) monomials, so by the lemma above, \(S\) must be a parallelogram.

This proves:
\[
n=4,\quad |\operatorname{supp}(F)|=4 \;\Longrightarrow\; S \text{ is the vertex set of a parallelogram}.
\]

---

## Conclusion

Under the stated setup, the **intended extremal classification** is:

\[
n=3 \implies |\operatorname{supp}(F)|\ge 4,
\]
so triangles cannot realize equality, and

\[
n=4,\ |\operatorname{supp}(F)|=4 \implies S \text{ is a parallelogram}.
\]

The literal intermediate claim
\[
|\operatorname{vert}(P+Q)|\ge 4
\]
in part \( (1) \) is false without an extra hypothesis, but the extremal conclusion itself is correct.